\section{Minimizing Coins}
\label{sec:cses-minimizing-coins}

\problemstatement{Given $n$ coin values and a target sum $x$, find the
\emph{minimum} number of coins needed to make exactly $x$ (each coin may
be used any number of times), or report $-1$ if it is impossible.}

\begin{patternbox}
\emph{``Minimum coins to make a sum,'' unlimited supply}, is the
\textbf{unbounded knapsack} minimisation. Let $\textit{dp}[s]$ be the
fewest coins summing to $s$; the last coin used is some value $c$, so
$\textit{dp}[s]=1+\min_c \textit{dp}[s-c]$.
\end{patternbox}

\subsection*{State and transition}

$\textit{dp}[s]$ = minimum coins to make sum $s$, with
$\textit{dp}[0]=0$ and all others initialised to ``infinity''
(unreachable). For each $s$, try every coin $c\le s$ as the last coin:
if $s-c$ is reachable, $s$ can be reached with one more coin. Take the
best:
\[
  \textit{dp}[s]=\min_{c:\,c\le s}\big(\textit{dp}[s-c]+1\big).
\]

\begin{ideabox}[Minimisation Instead of Counting]
This is Dice Combinations with $\min+1$ replacing summation: instead of
counting all ways to make the last move, we keep only the cheapest.
Coins may repeat, so every coin is available at every sum -- that is what
makes it \emph{unbounded}.
\end{ideabox}

\subsection*{Algorithm}

\begin{enumerate}
  \item $\textit{dp}[0]=0$; $\textit{dp}[s]=\infty$ for $s\ge1$.
  \item For $s=1\ldots x$ and each coin $c\le s$: if $\textit{dp}[s-c]$
        is finite, $\textit{dp}[s]=\min(\textit{dp}[s],
        \textit{dp}[s-c]+1)$.
  \item Answer is $\textit{dp}[x]$, or $-1$ if it is still $\infty$.
\end{enumerate}

\subsection*{Dry run}

\dryrun{Coins $\{1,5,7\}$, $x=11$.}

\begin{itemize}
  \item $\textit{dp}[5]=1$ (one 5), $\textit{dp}[7]=1$,
        $\textit{dp}[10]=2$ ($5+5$).
  \item $\textit{dp}[11]$: via coin $1$ from $\textit{dp}[10]=2\to3$;
        via coin $5$ from $\textit{dp}[6]$ ($=2$, as $5{+}1$)$\to3$; via
        coin $7$ from $\textit{dp}[4]$ ($=4$)$\to5$. Minimum $=3$.
  \item Answer $3$ (e.g.\ $5+5+1$).
\end{itemize}

\subsection*{C++ implementation}

\begin{lstlisting}
#include <bits/stdc++.h>
using namespace std;

int main() {
    int n, x;
    cin >> n >> x;
    vector<int> coins(n);
    for (int& c : coins) cin >> c;

    const int INF = 1e9;
    vector<int> dp(x + 1, INF);
    dp[0] = 0;
    for (int s = 1; s <= x; s++)
        for (int c : coins)
            if (c <= s && dp[s - c] != INF)
                dp[s] = min(dp[s], dp[s - c] + 1);

    cout << (dp[x] == INF ? -1 : dp[x]) << "\n";
    return 0;
}
\end{lstlisting}

\begin{complexitybox}
\textbf{Time Complexity.} $O(nx)$ -- each of $x$ sums tries all $n$
coins. \textbf{Space Complexity.} $O(x)$ for the DP array.
\end{complexitybox}

\begin{takeawaybox}
Swap ``sum of ways'' for ``$\min$ of ways $+1$'' and a counting DP becomes
an optimisation DP. The unbounded (coins reusable) structure is why every
coin is tried at every sum, distinguishing this from the $0/1$ knapsack of
Book Shop (Section~\ref{sec:cses-book-shop}).
\end{takeawaybox}
